% GENERATED FILE. Edit canonical lessons, then run npm run generate:books.
% canonical-source-hash: fnv1a64:710aa1a3a42353eb
% canonical-lessons: SA-C533-antarayah, SA-C533-avarodhah, SA-C533-stambhayati, SA-C533-apasarayati, SA-C533-kastam, SA-R533-review-526-533

\chapter{Something in the Way}
\label{ch:sa-something-in-the-way}
\hlchaptermodality{533}

\begin{chapteropening}
\textbf{By the end of this chapter:} \emph{I can say what will get in the way and what we can do.}

Complete the last lesson of chapter 533: I can say what will get in the way and what we can do.
\end{chapteropening}

\section[\texorpdfstring{\sk{अन्तरायः} (antarāyaḥ)}{अन्तरायः (antarāyaḥ)}]{\sk{अन्तरायः} (antarāyaḥ) — an impediment}

\label{lesson:SA-C533-antarayah}



\begin{quote}
Before the new one: say the Sanskrit for a girdle, then the Sanskrit for presents.
\end{quote}

\subsection*{You'll want to know: \sk{अन्तरायः}}
\textbf{\sk{अन्तरायः}} — \emph{antarāyaḥ} — \textquotedblleft{}an impediment\textquotedblright{}.

\begin{quote}
\sk{यदि} \sk{मार्गे} \sk{अन्तरायः} \sk{अस्ति}, \sk{तर्हि} \sk{वयम्} \sk{अन्यं} \sk{मार्गं} \sk{गच्छामः।}
\end{quote}

\begin{quote}
\emph{yadi mārge antarāyaḥ asti, tarhi vayam anyaṃ mārgaṃ gacchāmaḥ.}
\end{quote}

\begin{quote}
If there is something in the way, we take another road.
\end{quote}

A masculine noun.

\begin{cousinweb}
\textbf{\sk{अन्तरायः}} (\emph{antarāyaḥ}) is built on \emph{antar}, \textquotedblleft{}between\textquotedblright{}, and \emph{aya}, \textquotedblleft{}a going\textquotedblright{}, from the root √i, \textquotedblleft{}to go\textquotedblright{}: what comes between; MW gives \emph{antarāya} as intervention, an obstacle (MW).
\end{cousinweb}

\begin{grammarlens}[title={What you've built}]
One more word for saying what happens if.
\end{grammarlens}

\subsection*{Your turn}
\begin{itemize}
\raggedright
  \item \emph{Say it:} \emph{antarāyaḥ}
  \item \emph{Say it:} \emph{antarāyaḥ}, once more
  \item \emph{Say it:} \emph{yadi mārge antarāyaḥ asti, tarhi vayam anyaṃ mārgaṃ gacchāmaḥ.} — If there is something in the way, we take another road.
  \item \emph{Recall:} say the Sanskrit for a girdle, then the Sanskrit for presents, then say \emph{antarāyaḥ} again
\end{itemize}

\begin{tcolorbox}[breakable,colback=teal!4,colframe=teal!35!black,title={Before you move on}]
What does \sk{अन्तरायः} mean? (\textquotedblleft{}An impediment\textquotedblright{}.) Say it once more.
\end{tcolorbox}
\section[\texorpdfstring{\sk{अवरोधः} (avarodhaḥ)}{अवरोधः (avarodhaḥ)}]{\sk{अवरोधः} (avarodhaḥ) — an obstruction}

\label{lesson:SA-C533-avarodhah}



\begin{quote}
Before the new one: say the Sanskrit for presents, then the Sanskrit for an impediment.
\end{quote}

\subsection*{You'll want to know: \sk{अवरोधः}}
\textbf{\sk{अवरोधः}} — \emph{avarodhaḥ} — \textquotedblleft{}an obstruction\textquotedblright{}.

\begin{quote}
\sk{यदि} \sk{नद्याम्} \sk{अवरोधः} \sk{भवति}, \sk{तर्हि} \sk{जलं} \sk{न} \sk{प्रवहति।}
\end{quote}

\begin{quote}
\emph{yadi nadyām avarodhaḥ bhavati, tarhi jalaṃ na pravahati.}
\end{quote}

\begin{quote}
If the river is blocked, the water does not flow.
\end{quote}

A masculine noun.

\begin{cousinweb}
\textbf{\sk{अवरोधः}} (\emph{avarodhaḥ}) is built on \emph{ava}, \textquotedblleft{}down\textquotedblright{}, and the root √rudh, \textquotedblleft{}to obstruct\textquotedblright{}; MW gives \emph{avarodha} as hindrance, obstruction (MW).
\end{cousinweb}

\begin{grammarlens}[title={What you've built}]
One more word for saying what happens if.
\end{grammarlens}

\subsection*{Your turn}
\begin{itemize}
\raggedright
  \item \emph{Say it:} \emph{avarodhaḥ}
  \item \emph{Say it:} \emph{avarodhaḥ}, once more
  \item \emph{Say it:} \emph{yadi nadyām avarodhaḥ bhavati, tarhi jalaṃ na pravahati.} — If the river is blocked, the water does not flow.
  \item \emph{Recall:} say the Sanskrit for presents, then the Sanskrit for an impediment, then say \emph{avarodhaḥ} again
\end{itemize}

\begin{tcolorbox}[breakable,colback=teal!4,colframe=teal!35!black,title={Before you move on}]
What does \sk{अवरोधः} mean? (\textquotedblleft{}An obstruction\textquotedblright{}.) Say it once more.
\end{tcolorbox}
\section[\texorpdfstring{\sk{स्तम्भयति} (stambhayati)}{स्तम्भयति (stambhayati)}]{\sk{स्तम्भयति} (stambhayati) — brings to a standstill}

\label{lesson:SA-C533-stambhayati}



\begin{quote}
Before the new one: say the Sanskrit for an impediment, then the Sanskrit for an obstruction.
\end{quote}

\subsection*{You'll want to know: \sk{स्तम्भयति}}
\textbf{\sk{स्तम्भयति}} — \emph{stambhayati} — \textquotedblleft{}brings to a standstill\textquotedblright{}.

\begin{quote}
\sk{यदि} \sk{वर्षा} \sk{यानानि} \sk{स्तम्भयति}, \sk{तर्हि} \sk{वयं} \sk{श्वः} \sk{आगच्छामः।}
\end{quote}

\begin{quote}
\emph{yadi varṣā yānāni stambhayati, tarhi vayaṃ śvaḥ āgacchāmaḥ.}
\end{quote}

\begin{quote}
If the rain brings the traffic to a standstill, we will come tomorrow.
\end{quote}

Past (it …): \sk{वर्षा} \sk{यानानि} \sk{स्तम्भितवती}, \emph{varṣā yānāni stambhitavatī}.

\begin{cousinweb}
\textbf{\sk{स्तम्भयति}} (\emph{stambhayati}) comes from the causative of the root √stambh, \textquotedblleft{}to prop, to make firm\textquotedblright{}, the root of \emph{stambha}, \textquotedblleft{}a pillar\textquotedblright{}; MW gives \emph{stambhayati} as to fix, to make stiff, and to stop, to check (MW).
\end{cousinweb}

\begin{grammarlens}[title={What you've built}]
One more word for saying what happens if.
\end{grammarlens}

\subsection*{Your turn}
\begin{itemize}
\raggedright
  \item \emph{Say it:} \emph{stambhayati}
  \item \emph{Say it:} \emph{stambhayati}, once more
  \item \emph{Say it:} \emph{yadi varṣā yānāni stambhayati, tarhi vayaṃ śvaḥ āgacchāmaḥ.} — If the rain brings the traffic to a standstill, we will come tomorrow.
  \item \emph{Recall:} say the Sanskrit for an impediment, then the Sanskrit for an obstruction, then say \emph{stambhayati} again
\end{itemize}

\begin{tcolorbox}[breakable,colback=teal!4,colframe=teal!35!black,title={Before you move on}]
What does \sk{स्तम्भयति} mean? (\textquotedblleft{}Brings to a standstill\textquotedblright{}.) Say it once more.
\end{tcolorbox}
\section[\texorpdfstring{\sk{अपसारयति} (apasārayati)}{अपसारयति (apasārayati)}]{\sk{अपसारयति} (apasārayati) — moves out of the way}

\label{lesson:SA-C533-apasarayati}



\begin{quote}
Before the new one: say the Sanskrit for an obstruction, then the Sanskrit for brings to a standstill.
\end{quote}

\subsection*{You'll want to know: \sk{अपसारयति}}
\textbf{\sk{अपसारयति}} — \emph{apasārayati} — \textquotedblleft{}moves out of the way\textquotedblright{}.

\begin{quote}
\sk{यदि} \sk{मार्गे} \sk{शिला} \sk{अस्ति}, \sk{तर्हि} \sk{सः} \sk{ताम्} \sk{अपसारयति।}
\end{quote}

\begin{quote}
\emph{yadi mārge śilā asti, tarhi saḥ tām apasārayati.}
\end{quote}

\begin{quote}
If there is a stone on the road, he moves it out of the way.
\end{quote}

Past (he …): \sk{सः} \sk{शिलाम्} \sk{अपसारितवान्}, \emph{saḥ śilām apasāritavān}.

\begin{cousinweb}
\textbf{\sk{अपसारयति}} (\emph{apasārayati}) is built on \emph{apa}, \textquotedblleft{}away\textquotedblright{}, and the causative of the root √sṛ, \textquotedblleft{}to run, to flow\textquotedblright{}; MW gives \emph{apasārayati} as to make or let go away, to remove (MW).
\end{cousinweb}

\begin{grammarlens}[title={What you've built}]
One more word for saying what happens if.
\end{grammarlens}

\subsection*{Your turn}
\begin{itemize}
\raggedright
  \item \emph{Say it:} \emph{apasārayati}
  \item \emph{Say it:} \emph{apasārayati}, once more
  \item \emph{Say it:} \emph{yadi mārge śilā asti, tarhi saḥ tām apasārayati.} — If there is a stone on the road, he moves it out of the way.
  \item \emph{Recall:} say the Sanskrit for an obstruction, then the Sanskrit for brings to a standstill, then say \emph{apasārayati} again
\end{itemize}

\begin{tcolorbox}[breakable,colback=teal!4,colframe=teal!35!black,title={Before you move on}]
What does \sk{अपसारयति} mean? (\textquotedblleft{}Moves out of the way\textquotedblright{}.) Say it once more.
\end{tcolorbox}
\section[\texorpdfstring{\sk{कष्टम्} (kaṣṭam)}{कष्टम् (kaṣṭam)}]{\sk{कष्टम्} (kaṣṭam) — trouble; hardship}

\label{lesson:SA-C533-kastam}



\begin{quote}
Before the new one: say the Sanskrit for brings to a standstill, then the Sanskrit for moves out of the way.
\end{quote}

\subsection*{You'll want to know: \sk{कष्टम्}}
\textbf{\sk{कष्टम्}} — \emph{kaṣṭam} — \textquotedblleft{}trouble; hardship\textquotedblright{}.

\begin{quote}
\sk{यदि} \sk{कष्टम्} \sk{आगच्छति}, \sk{तर्हि} \sk{धैर्यं} \sk{धारयतु।}
\end{quote}

\begin{quote}
\emph{yadi kaṣṭam āgacchati, tarhi dhairyaṃ dhārayatu.}
\end{quote}

\begin{quote}
If trouble comes, hold on to your courage.
\end{quote}

A neuter noun.

\begin{cousinweb}
\textbf{\sk{कष्टम्}} (\emph{kaṣṭam}) is the word MW gives as bad, ill, wrong, and as a noun evil, trouble; MW suggests, with a \textquotedblleft{}perhaps\textquotedblright{}, the past participle of the root √kaṣ, \textquotedblleft{}to rub, to scratch\textquotedblright{}, so its origin is uncertain (MW).
\end{cousinweb}

\begin{grammarlens}[title={What you've built}]
One more word for saying what happens if.
\end{grammarlens}

\subsection*{Your turn}
\begin{itemize}
\raggedright
  \item \emph{Say it:} \emph{kaṣṭam}
  \item \emph{Say it:} \emph{kaṣṭam}, once more
  \item \emph{Say it:} \emph{yadi kaṣṭam āgacchati, tarhi dhairyaṃ dhārayatu.} — If trouble comes, hold on to your courage.
  \item \emph{Recall:} say the Sanskrit for brings to a standstill, then the Sanskrit for moves out of the way, then say \emph{kaṣṭam} again
\end{itemize}

\begin{tcolorbox}[breakable,colback=teal!4,colframe=teal!35!black,title={Before you move on}]
What does \sk{कष्टम्} mean? (\textquotedblleft{}Trouble; hardship\textquotedblright{}.) Say it once more.
\end{tcolorbox}
\section[(dialogue)]{Looking back: from Training Hard to Something in the Way}

\label{lesson:SA-R533-review-526-533}



\begin{quote}
Say the Sanskrit for a gymnasium and a runner.
\end{quote}

\subsection*{Your turn}
Say these aloud:

\begin{itemize}
\raggedright
  \item a gymnasium, a runner, quick, leaps over, a team
  \item an archer, focused, shoots, cheers on, a bowstring
  \item organises, materials, supplies, makes ready, an almanac
  \item a change, winds up, revises, hurries along, haste
  \item eatables, a drink, a laddu, a sherbet, stale
  \item yearns for, a heart's wish, propitiates, a prayer, a desire
  \item a present, a bunch of flowers, fine cloth, a girdle, presents
  \item an impediment, an obstruction, brings to a standstill, moves out of the way, trouble
\end{itemize}

\begin{tcolorbox}[breakable,colback=teal!4,colframe=teal!35!black,title={Before you move on}]
A gymnasium? (\textbf{\sk{व्यायामभूमिः}}, \emph{vyāyāmabhūmiḥ}.) A runner? (\textbf{\sk{धावकः}}, \emph{dhāvakaḥ}.) Quick? (\textbf{\sk{क्षिप्रः}}, \emph{kṣipraḥ}.) Leaps over? (\textbf{\sk{लंघयति}}, \emph{laṅghayati}.) A team? (\textbf{\sk{दलम्}}, \emph{dalam}.) An archer? (\textbf{\sk{धनुर्धरः}}, \emph{dhanurdharaḥ}.) Focused? (\textbf{\sk{एकाग्रः}}, \emph{ekāgraḥ}.) Shoots? (\textbf{\sk{प्रहिणोति}}, \emph{prahiṇoti}.) Cheers on? (\textbf{\sk{उत्साहयति}}, \emph{utsāhayati}.) A bowstring? (\textbf{\sk{ज्या}}, \emph{jyā}.) Organises? (\textbf{\sk{आयोजयति}}, \emph{āyojayati}.) Materials? (\textbf{\sk{सामग्री}}, \emph{sāmagrī}.) Supplies? (\textbf{\sk{सम्भारः}}, \emph{sambhāraḥ}.) Makes ready? (\textbf{\sk{उपकल्पयति}}, \emph{upakalpayati}.) An almanac? (\textbf{\sk{पञ्चांगम्}}, \emph{pañcāṅgam}.) A change? (\textbf{\sk{परिवर्तनम्}}, \emph{parivartanam}.) Winds up? (\textbf{\sk{उपसंहरति}}, \emph{upasaṃharati}.) Revises? (\textbf{\sk{संशोधयति}}, \emph{saṃśodhayati}.) Hurries along? (\textbf{\sk{त्वरयति}}, \emph{tvarayati}.) Haste? (\textbf{\sk{त्वरा}}, \emph{tvarā}.) Eatables? (\textbf{\sk{भोज्यम्}}, \emph{bhojyam}.) A drink? (\textbf{\sk{पेयम्}}, \emph{peyam}.) A laddu? (\textbf{\sk{लड्डुकः}}, \emph{laḍḍukaḥ}.) A sherbet? (\textbf{\sk{पानकम्}}, \emph{pānakam}.) Stale? (\textbf{\sk{पर्युषितः}}, \emph{paryuṣitaḥ}.) Yearns for? (\textbf{\sk{वाञ्छति}}, \emph{vāñchati}.) A heart's wish? (\textbf{\sk{मनोरथः}}, \emph{manorathaḥ}.) Propitiates? (\textbf{\sk{आराधयति}}, \emph{ārādhayati}.) A prayer? (\textbf{\sk{प्रार्थना}}, \emph{prārthanā}.) A desire? (\textbf{\sk{कामना}}, \emph{kāmanā}.) A present? (\textbf{\sk{उपायनम्}}, \emph{upāyanam}.) A bunch of flowers? (\textbf{\sk{गुच्छः}}, \emph{gucchaḥ}.) Fine cloth? (\textbf{\sk{दुकूलम्}}, \emph{dukūlam}.) A girdle? (\textbf{\sk{मेखला}}, \emph{mekhalā}.) Presents? (\textbf{\sk{उपहरति}}, \emph{upaharati}.) An impediment? (\textbf{\sk{अन्तरायः}}, \emph{antarāyaḥ}.) An obstruction? (\textbf{\sk{अवरोधः}}, \emph{avarodhaḥ}.) Brings to a standstill? (\textbf{\sk{स्तम्भयति}}, \emph{stambhayati}.) Moves out of the way? (\textbf{\sk{अपसारयति}}, \emph{apasārayati}.) Trouble? (\textbf{\sk{कष्टम्}}, \emph{kaṣṭam}.)
\end{tcolorbox}
